% Standalone problem document, generated from the adjacent README.md.
% Compile with XeLaTeX. Mathematical content is maintained in Markdown.
\documentclass[11pt,a4paper]{article}
\usepackage[fontsize=10.5pt]{scrextend}
\usepackage[margin=22mm,headheight=15pt,headsep=9mm,footskip=11mm]{geometry}
\usepackage{fontspec}
\setmainfont{texgyrepagella}[Extension=.otf,UprightFont=*-regular,BoldFont=*-bold,ItalicFont=*-italic,BoldItalicFont=*-bolditalic]
\setsansfont{texgyreheros}[Extension=.otf,UprightFont=*-regular,BoldFont=*-bold,ItalicFont=*-italic,BoldItalicFont=*-bolditalic]
\setmonofont{texgyrecursor}[Extension=.otf,UprightFont=*-regular,BoldFont=*-bold,ItalicFont=*-italic,BoldItalicFont=*-bolditalic,Scale=MatchLowercase]
\usepackage{amsmath,amssymb}
\usepackage{unicode-math}
\setmathfont{texgyrepagella-math.otf}
\usepackage{xcolor,microtype,enumitem,needspace,fancyhdr}
\usepackage{bookmark}
\definecolor{ink}{HTML}{17344B}
\definecolor{accent}{HTML}{197D80}
\definecolor{muted}{HTML}{596773}
\hypersetup{colorlinks=true,linkcolor=accent,urlcolor=accent,pdftitle={IV-06 - Number
of components of a real interval eigenvalue
set},pdfauthor={Open Problems in Numerical Linear Algebra}}
\urlstyle{same}
\setlength{\parindent}{0pt}
\setlength{\parskip}{5pt plus 1pt minus 1pt}
\linespread{1.04}
\setlength{\emergencystretch}{3em}
\widowpenalty=10000
\clubpenalty=10000
\displaywidowpenalty=10000
\allowdisplaybreaks[1]
\setlist{nosep,leftmargin=1.5em}
\providecommand{\tightlist}{\setlength{\itemsep}{0pt}\setlength{\parskip}{0pt}}
\providecommand{\pandocbounded}[1]{#1}
\pagestyle{fancy}
\fancyhf{}
\fancyhead[L]{\sffamily\footnotesize\color{muted}OPEN PROBLEMS IN NUMERICAL LINEAR ALGEBRA}
\fancyhead[R]{\sffamily\bfseries\footnotesize\color{accent}IV-06}
\fancyfoot[L]{\parbox[b]{0.9\textwidth}{\sffamily\fontsize{8}{10}\selectfont\color{muted}Editorial ratings; a bounded search cannot certify that a problem remains unsolved.\\Literature check: 2026-09-11}}
\fancyfoot[R]{\sffamily\footnotesize\color{muted}\thepage}
\renewcommand{\headrulewidth}{0.3pt}
\renewcommand{\headrule}{\hbox to\headwidth{\color{accent}\leaders\hrule height \headrulewidth\hfill}}
\makeatletter
\renewcommand{\section}{\Needspace{5\baselineskip}\@startsection{section}{1}{0pt}{12pt plus 2pt minus 2pt}{4pt}{\sffamily\bfseries\large\color{ink}}}
\renewcommand{\subsection}{\Needspace{4\baselineskip}\@startsection{subsection}{2}{0pt}{9pt plus 1pt minus 1pt}{3pt}{\sffamily\bfseries\normalsize\color{ink}}}
\makeatother
\setcounter{secnumdepth}{0}
\begin{document}
{\sffamily\small\color{muted}Intervals and absolute value equations\par}
\vspace{3pt}
{\raggedright\hyphenpenalty=10000\exhyphenpenalty=10000\sffamily\bfseries\fontsize{21}{24}\selectfont\color{ink}Number
of components of a real interval eigenvalue set\par}
\vspace{7pt}
{\sffamily\small\textbf{Difficulty:} challenging\quad\textbf{Importance:} interesting
to specialist\par}
{\sffamily\small\bfseries\color{accent}Status: Lean verified\par}
\vspace{4pt}
{\color{accent}\hrule height 0.8pt}
\vspace{6pt}
\textbf{Rating rationale:} Controlling the topology of the real spectrum
for nonsymmetric uncertain matrices appears to require new analysis,
warranting challenging. The component-count bound primarily advances
specialist interval spectral enclosure theory.

\textbf{Area:} interval linear algebra; eigenvalue computation

\subsection{Independently reviewed resolution -
2026-09-11}\label{independently-reviewed-resolution---2026-09-11}

\textbf{Negative resolution.} Theorem 1 gives a \(3\times3\)
independent-entry interval matrix with at least four components in its
real eigenvalue set. Four exact integer eigenpairs at \(-3,0,3,25\) and
excluded separators \(-1,1,12\) refute the universal at-most-\(n\)
conjecture. No symmetry assumption is introduced, and locating every
component endpoint is unnecessary.

\textbf{Author:} Matthew J. Colbrook, Department of Applied Mathematics
and Theoretical Physics, University of Cambridge. See the
\href{https://github.com/ajt60gaibb/OpenProblemsInNLA/blob/main/references/colbrook-intervals-2026-09-11/manuscripts/IV-06.pdf}{complete
manuscript},
\href{https://github.com/ajt60gaibb/OpenProblemsInNLA/blob/main/references/colbrook-intervals-2026-09-11/verification/reviews/IV-06-review.md}{independent
agent review} and
\href{https://github.com/ajt60gaibb/OpenProblemsInNLA/blob/main/references/colbrook-intervals-2026-09-11/README.md}{submission
record}. The source archive identifies the drafts as AI-generated;
authorship is recorded at the submitter's request. This is agent
verification, not external human peer review or formal proof-assistant
certification.

The difficulty, importance and rating rationale below are historical
assessments of the original open target. The original statement and
dated audits are preserved.

\subsection{Lean verification ---
2026-09-14}\label{lean-verification-2026-09-14}

\textbf{The complete original conjecture has a Lean-verified negative
answer.}
\href{https://github.com/sidneyholden1/OpenProblemsInNLA/tree/23dd1e1727494efac35743826c0725ebc71942c6/intervals-and-absolute-value-equations/IV-06/lean}{The
formal proof} uses actual nonzero real eigenvectors and actual connected
components of the real eigenvalue set. Four exact eigenpairs and three
excluded separators give at least four components in dimension three.
The topological component-count argument is formally proved.

\textbf{Formalization:} Sidney Holden, Center for Computational Biology,
Flatiron Institute, Simons Foundation, with OpenAI Codex assistance.
Matthew J. Colbrook retains authorship of the mathematical
counterexample. Two independent AI statement reviews preceded
implementation, and two independent final proof reviews approved the
complete formalization. Kernel-only LeanCert, all four exported axiom
audits, and
\href{https://github.com/sidneyholden1/OpenProblemsInNLA/actions/runs/34805467886}{isolated
Linux Comparator and kernel checks} passed at the immutable revision
linked above.
\href{https://github.com/ajt60gaibb/OpenProblemsInNLA/blob/main/intervals-and-absolute-value-equations/IV-06/lean/verification/linux-2026-09-14/OPERATIONAL-REVIEW.md}{Original
evidence and operational audit} retain the exact 26 checked inputs and
original logs. The first Linux attempt failed during a checker download;
the unchanged proof passed on retry. No external human peer review or
source-author endorsement is asserted.

\newpage

\subsection{Problem statement}\label{problem-statement}

For every integer \(n\ge1\) and real matrices
\(L,U\in\mathbb R^{n\times n}\) with \(L\le U\) entrywise, form the
interval matrix family

\[\mathcal A=\{A\in\mathbb R^{n\times n}:L\le A\le U\}.\]

Each entry varies independently in its prescribed closed interval;
singleton intervals are allowed. Define the set of all real eigenvalues
attained by the family by

\[\Lambda_{\mathbb R}(\mathcal A)
=\{\lambda\in\mathbb R:\exists A\in\mathcal A,\ \exists x\in\mathbb R^n\setminus\{0\},\ Ax=\lambda x\}.\]

Must \(\Lambda_{\mathbb R}(\mathcal A)\) have at most \(n\) connected
components? Equivalently, must there exist \(0\le m\le n\) and real
numbers \(a_r\le b_r\), \(1\le r\le m\), such that

\[\Lambda_{\mathbb R}(\mathcal A)=\bigcup_{r=1}^{m}[a_r,b_r]?\]

The empty set corresponds to \(m=0\), and a singleton eigenvalue
component counts as one interval. No symmetry, diagonalizability, or
reality of the entire spectrum is assumed. Only real eigenvalues
contribute to the set.

Interval eigensolvers enclose all eigenvalues allowed by uncertain
matrix entries. The conjecture bounds the number of separated real
spectral regions they must represent, without claiming that their
endpoints can be found efficiently. Finiteness and compactness of the
union are known; the proposed linear bound is the unresolved claim. For
families restricted to symmetric matrices, the ordered eigenvalue
functions already give \(n\) compact intervals, as explained in the
second reference; that does not address the family above.

\newpage
\subsection{References}\label{references}

Milan Hladík, David Daney, and Elias P. Tsigaridas,
\href{https://doi.org/10.1016/j.cam.2010.11.022}{\emph{An algorithm for
addressing the real interval eigenvalue problem}}, Journal of
Computational and Applied Mathematics \textbf{235} (2011), 2715--2730,
§1, p.~2716, paragraph immediately after the definition of \(\Lambda\).

The same authors,
\href{https://www-sop.inria.fr/coprin/PDF/CAMWA6430.pdf}{\emph{Characterizing
and approximating eigenvalue sets of symmetric interval matrices}},
Computers \& Mathematics with Applications \textbf{62} (2011),
3152--3163, §2, p.~3153;
\href{https://doi.org/10.1016/j.camwa.2011.08.028}{DOI}.

\subsection{Status check}\label{status-check}

On 2026-09-10, checked the published formulation, the authors' follow-up
on symmetric interval matrices, Hladík's current publication record, and
searches for the original title, interval eigenvalue components, the
at-most-\(n\) claim, and 2025/2026 developments. No primary source
proving or refuting the general conjecture was located. This is a
bounded literature check; papers addressing only symmetric interval
families do not settle the stated question.

\subsection{Independent audit ---
2026-09-10}\label{independent-audit-2026-09-10}

The primary 2011 JCAM text explicitly poses the at-most-\(n\) conjecture
for independent entries. The full symmetric follow-up instead studies
the coupled family \(\mathcal A^S=\{A\in\mathcal A:A=A^T\}\), generally
a proper subset, as its §2 explains. Its result alone does not justify
the previous partial-resolution label for the displayed
independent-entry target. Targeted later component-count and
exact-conjecture searches found no proof or counterexample; the explicit
openness statement is historical.
\end{document}
